% Master's thesis — classical vs quantum vs hybrid machine learning.
% Chapter map and page budgets: docs/groundwork.md §2.
%
% TEMPLATE NOTE: this preamble is a deliberately plain, template-agnostic base
% (university template = open question 1 in groundwork §12). When the official
% class arrives, replace this preamble; the chapter files are plain LaTeX and
% port unchanged.
\documentclass[11pt,a4paper,oneside]{report}

\usepackage[utf8]{inputenc}
\usepackage[T1]{fontenc}
\usepackage{lmodern}
\usepackage{microtype}
\usepackage[margin=2.6cm]{geometry}
\usepackage{amsmath,amssymb,amsthm,bm}
\usepackage{graphicx}
\usepackage{booktabs}
\usepackage{xcolor}
\usepackage[numbers,sort&compress]{natbib}
\usepackage[hidelinks]{hyperref}

% Shared notation — keep ALL notation choices here so chapters stay consistent.

% Dirac notation (kept self-contained; no `physics` package for portability)
\newcommand{\ket}[1]{\left|#1\right\rangle}
\newcommand{\bra}[1]{\left\langle#1\right|}
\newcommand{\braket}[2]{\left\langle#1\middle|#2\right\rangle}
\newcommand{\ketbra}[2]{\left|#1\middle\rangle\middle\langle#2\right|}

% Operators and sets
\newcommand{\E}{\mathbb{E}}
\newcommand{\Var}{\mathrm{Var}}
\newcommand{\Tr}{\mathrm{Tr}}
\newcommand{\R}{\mathbb{R}}
\newcommand{\C}{\mathbb{C}}
\newcommand{\Hil}{\mathcal{H}}
\newcommand{\Dcal}{\mathcal{D}}
\newcommand{\Fcal}{\mathcal{F}}
\newcommand{\norm}[1]{\left\lVert#1\right\rVert}
\newcommand{\abs}[1]{\left\lvert#1\right\rvert}
\DeclareMathOperator{\sign}{sign}
\DeclareMathOperator*{\argmin}{arg\,min}
\DeclareMathOperator*{\argmax}{arg\,max}

% Recurring thesis-specific symbols
\newcommand{\data}{\mathcal{S}}            % dataset / sample
\newcommand{\enc}[1]{S(#1)}                % data-encoding circuit
\newcommand{\ansatz}{W(\bm{\theta})}       % trainable ansatz
\newcommand{\qkernel}{\kappa_q}            % quantum (fidelity) kernel

% Theorem-like environments
\theoremstyle{definition}
\newtheorem{definition}{Definition}[chapter]
\newtheorem{example}{Example}[chapter]
\theoremstyle{plain}
\newtheorem{theorem}{Theorem}[chapter]
\newtheorem{proposition}{Proposition}[chapter]
\theoremstyle{remark}
\newtheorem{remark}{Remark}[chapter]

% Draft annotations — must all be resolved before submission
\newcommand{\todonote}[1]{\textcolor{orange!80!black}{[TODO: #1]}}
\newcommand{\resultref}[1]{\textcolor{teal!70!black}{[RESULT: #1]}}  % filled from analysis/ scripts only


\title{Classical, Quantum, and Hybrid Machine Learning:\\
       A Systematic Empirical Comparison under Matched Conditions}
\author{\todonote{Author name}\\ \todonote{Institution, programme}}
\date{\todonote{Submission date}}

\begin{document}

\maketitle

\begin{abstract}
\todonote{Abstract is written last (week 19). One paragraph each: context and
the hype--evidence tension; the five research questions; the matched-conditions
methodology (shared splits, matched tuning budgets, parameter matching,
regime-separated evaluation); headline empirical findings; the honest-framing
conclusion about NISQ-scale competitiveness rather than quantum advantage.}
\end{abstract}

\tableofcontents

\chapter{Introduction}
\label{ch:introduction}

\section{Motivation}
\label{sec:intro-motivation}

Machine learning and quantum computing are, individually, two of the most
consequential computational developments of the past two decades. Their
intersection --- \emph{quantum machine learning} (QML) --- has attracted
extraordinary attention, driven by the prospect that quantum computers might
learn from data in ways that classical computers cannot
\citep{biamonte2017qml, schuld2021book}. Parameterized quantum circuits can
represent functions in Hilbert spaces whose dimension grows exponentially with
the number of qubits, and quantum kernels can measure similarities between data
points embedded in those spaces \citep{havlicek2019kernels, schuld2019hilbert}.
On paper, this looks like a recipe for more expressive models.

The empirical record is far more ambiguous. The devices available today ---
noisy intermediate-scale quantum (NISQ) hardware \citep{preskill2018nisq} ---
support only small, shallow circuits. Provable quantum--classical learning
separations exist but are confined to carefully constructed problems
\citep{liu2021rigorous}, and a growing body of work shows that classical models
given access to the same data often match or exceed quantum models on the
benchmarks actually used in the literature \citep{huang2021power,
bowles2024better, schnabel2025scrutiny}. At the same time, methodological
audits of published QML studies find recurring flaws that systematically favor
the quantum contender: tuned quantum models compared against default-configured
classical baselines, cherry-picked datasets, inconsistent data splits, and
metrics aggregated across incompatible execution regimes
\citep{bowles2024better}.

This thesis takes the tension between promise and evidence as its starting
point. Rather than asking whether quantum machine learning is ``better'' in the
abstract, it asks a narrower, answerable question: \emph{under carefully matched
conditions --- identical data splits, identical preprocessing, identical
hyperparameter-search budgets, and parameter-matched architectures --- how do
quantum and hybrid models actually compare to strong classical baselines on
standard benchmark tasks, and which design factors drive the differences?}
Answering this requires as much care in experimental methodology as in quantum
theory, and the methodological framework is itself one of the main
contributions of this work.

\section{Research questions}
\label{sec:intro-rqs}

Five research questions drive both the theoretical treatment and the
experimental programme of this thesis.

\begin{description}
  \item[RQ1 (Performance).] Under matched data, preprocessing, and
  hyperparameter-search budgets, how do quantum models --- quantum kernels and
  variational quantum classifiers --- compare to strong classical baselines in
  accuracy, balanced accuracy, $F_1$, and ROC-AUC?

  \item[RQ2 (Design factors).] How do the data-encoding strategy, ansatz
  choice, circuit depth, and entanglement structure affect quantum-model
  performance and trainability?

  \item[RQ3 (Hybrid value).] Do hybrid architectures --- classical feature
  extraction feeding a quantum head, quantum kernels inside a classical SVM,
  quantum transfer learning --- outperform either pure approach at matched
  parameter counts?

  \item[RQ4 (Robustness and cost).] How do finite sampling (shot noise),
  simulated device noise, and real quantum hardware affect the results, and
  what are the actual resource costs --- wall-clock time, parameter counts,
  circuit depth, and shot budgets?

  \item[RQ5 (Data regime).] Does the relative performance of classical,
  quantum, and hybrid models change with the training-set size and the feature
  dimensionality?
\end{description}

\section{Scope and honest framing}
\label{sec:intro-framing}

A thesis-scale experimental programme can reach system sizes of roughly ten to
twenty qubits. Every quantum model studied here can therefore be simulated
classically --- indeed, most experiments in this thesis \emph{are} exact
statevector simulations, by design. It follows that \textbf{this thesis cannot,
even in principle, demonstrate quantum advantage}: any task its quantum models
solve is, by construction, solvable classically at the same scale. What it can
do --- and what the strongest recent benchmark literature argues is currently
the most valuable contribution \citep{bowles2024better, schnabel2025scrutiny}
--- is measure the \emph{empirical competitiveness and behavior of quantum
models at NISQ scale} under a protocol fair enough that the observed
differences mean something. Claims in this thesis are therefore claims about
performance, trainability, robustness, and cost at the scales studied, never
about asymptotic separation. Where the literature offers provable separations
\citep{liu2021rigorous} or arguments about the role of data access
\citep{huang2021power}, they are treated in Chapter~\ref{ch:related} as context
for interpreting the empirical results, not as claims this thesis tests at
scale.

The scope is further restricted to \emph{supervised binary classification} on
standard benchmark datasets --- synthetic constructions with controlled
structure, small tabular datasets, and reduced image data --- the setting in
which the comparative literature is concentrated and in which matched
conditions are achievable within a Master's project. Multi-class extensions,
regression, generative models, and learning on inherently quantum data are
discussed as outlook.

\section{Contributions}
\label{sec:intro-contributions}

The main contributions of this thesis are:

\begin{enumerate}
  \item \textbf{A fairness-first experimental protocol}, specified before any
  quantum model was trained: all models --- classical, quantum, and hybrid ---
  share identical stored cross-validation splits (5-fold stratified, five
  seeds), identical $[0,\pi]$ feature scaling fit on training folds only,
  an identical hyperparameter-optimization budget (one 50-trial Optuna study
  per model--dataset pair under a fixed procedure), parameter-matched
  neural-network baselines, and strict separation of execution regimes (exact
  simulation, finite shots, noisy simulation, hardware) in all reported
  results.

  \item \textbf{A systematic empirical comparison} of seven classical
  baselines, quantum kernel methods, variational quantum classifiers, and
  hybrid architectures across ten benchmark datasets spanning saturated to
  genuinely hard, including controlled synthetic tasks (interleaved half-moons,
  concentric circles, and a sign-parity task with distractor dimensions
  designed to separate model families).

  \item \textbf{Focused design-factor studies}: a data-encoding ablation
  connected to the Fourier-expressivity theory of parameterized circuits
  \citep{schuld2021fourier}; a trainability study reproducing barren-plateau
  scaling \citep{mcclean2018barren}; sample-efficiency learning curves; and a
  noise ladder from exact simulation through calibrated device-noise models to
  validation runs on IBM quantum hardware.

  \item \textbf{A fully reproducible, open research artifact}: every figure and
  table in this document is generated by a version-controlled script from
  append-only result logs carrying configuration hashes, seeds, and split
  identifiers; the stored splits and locked hyperparameter configurations are
  part of the repository.
\end{enumerate}

\section{Thesis outline}
\label{sec:intro-outline}

Chapter~\ref{ch:classical} reviews the classical machine-learning foundations
the comparison rests on, with emphasis on kernel methods and evaluation
methodology. Chapter~\ref{ch:quantum} introduces the quantum-computing
concepts required for the remainder: states, gates, measurement, and the
realities of NISQ devices. Chapter~\ref{ch:qml} develops quantum machine
learning in depth --- data encodings and their expressivity consequences,
variational quantum algorithms, quantum kernel methods, the main circuit
architectures, trainability obstacles, and hybrid designs.
Chapter~\ref{ch:related} surveys comparative studies and the quantum-advantage
debate, positioning this thesis within that literature.
Chapter~\ref{ch:methodology} specifies the datasets, the model zoo, and the
fairness protocol in full detail. Chapter~\ref{ch:experiments} reports the
experimental results, one experiment family at a time, each stated as a
hypothesis before its outcome. Chapter~\ref{ch:discussion} synthesizes the
findings across research questions and examines threats to validity, and
Chapter~\ref{ch:conclusion} concludes with directions for future work.

\chapter{Classical Machine Learning Foundations}
\label{ch:classical}

The comparison at the heart of this thesis is only as strong as its classical
side. This chapter fixes the supervised-learning formalism and notation, then
reviews the model families that serve as baselines --- with particular depth on
kernel methods, since the mathematics of Chapter~\ref{ch:qml}'s quantum kernels
is the same mathematics --- and closes with the evaluation methodology that the
fairness protocol of Chapter~\ref{ch:methodology} operationalizes. Standard
references for this material are \citet{bishop2006prml} and
\citet{hastie2009esl}.

\section{Supervised learning}
\label{sec:ml-setup}

In supervised binary classification we observe a sample
$\data = \{(\bm{x}_i, y_i)\}_{i=1}^{m}$ drawn i.i.d.\ from an unknown
distribution $\Dcal$ over $\mathcal{X} \times \mathcal{Y}$ with
$\mathcal{X} \subseteq \R^d$ and $\mathcal{Y} = \{0, 1\}$. A learning
algorithm selects a hypothesis $h \colon \mathcal{X} \to \mathcal{Y}$ (often
via a real-valued score $f$ with $h = \mathbf{1}[f \ge \tau]$) from a
hypothesis class $\Fcal$, seeking small \emph{risk}
\begin{equation}
  R(h) = \E_{(\bm{x}, y) \sim \Dcal} \bigl[ \ell(h(\bm{x}), y) \bigr],
\end{equation}
for a loss $\ell$, while only the \emph{empirical risk}
$\widehat{R}(h) = \tfrac{1}{m} \sum_i \ell(h(\bm{x}_i), y_i)$ is observable.
The gap between the two is the object of generalization theory; its practical
management --- model selection without contaminating the estimate of $R$ ---
is the subject of Section~\ref{sec:ml-evaluation}.

The classical decomposition of expected error into \emph{bias},
\emph{variance}, and irreducible noise \citep{hastie2009esl} shapes intuition
throughout this thesis: rich hypothesis classes (deep trees, wide networks,
expressive quantum circuits) reduce bias at the price of variance, and
regularization --- explicit penalties, restricted architectures, or early
stopping --- trades back along the same curve. Two caveats are important.
Modern over-parameterized models can defy the naive U-shaped picture
\citep{goodfellow2016deep}, and the same caution applies to quantum models:
expressivity arguments alone predict neither good nor bad generalization
\citep{gilfuster2024rethinking}, which is why this thesis measures rather than
assumes (RQ2, RQ5).

\section{Linear models and logistic regression}
\label{sec:ml-linear}

The simplest score is affine, $f(\bm{x}) = \bm{w}^{\top} \bm{x} + b$.
Logistic regression passes it through the sigmoid
$\sigma(t) = (1 + e^{-t})^{-1}$, interprets
$\sigma(f(\bm{x}))$ as $\Pr(y = 1 \mid \bm{x})$, and fits $(\bm{w}, b)$ by
minimizing the $\ell_2$-regularized negative log-likelihood
\begin{equation}
  \min_{\bm{w}, b} \; \frac{1}{C} \cdot \frac{\norm{\bm{w}}^2}{2}
  + \sum_{i=1}^{m} \log\!\left( 1 + e^{-\tilde{y}_i (\bm{w}^{\top}\bm{x}_i + b)} \right),
  \qquad \tilde{y}_i = 2y_i - 1 \in \{\pm 1\},
\end{equation}
a smooth convex problem. Logistic regression contributes two things to this
thesis: a floor that any nonlinear model must clear to justify its complexity,
and a diagnostic --- tasks on which it performs at chance (the concentric
circles and sign-parity constructions of Chapter~\ref{ch:methodology}) are
certified to require nonlinear structure.

\section{Kernel methods and support vector machines}
\label{sec:ml-kernels}

Kernel methods make nonlinearity affordable by moving it into an inner
product. A \emph{feature map} $\phi \colon \mathcal{X} \to \Hil$ into a
(possibly infinite-dimensional) Hilbert space defines the \emph{kernel}
\begin{equation}
  \kappa(\bm{x}, \bm{x}') = \braket{\phi(\bm{x})}{\phi(\bm{x}')}_{\Hil},
  \label{eq:kernel}
\end{equation}
which is symmetric and positive semidefinite; conversely (Mercer), any such
kernel is an inner product in \emph{some} feature space. Algorithms that touch
data only through inner products can substitute $\kappa$ and operate in
$\Hil$ implicitly --- the \emph{kernel trick}.

The support vector machine (SVM) \citep{cortes1995svm} seeks the
maximum-margin separator in $\Hil$. Its soft-margin primal,
\begin{equation}
  \min_{\bm{w}, b, \bm{\xi}} \;\; \frac{1}{2} \norm{\bm{w}}^2
    + C \sum_{i=1}^{m} \xi_i
  \quad \text{s.t.} \quad
  \tilde{y}_i \bigl( \braket{\bm{w}}{\phi(\bm{x}_i)} + b \bigr) \ge 1 - \xi_i,
  \;\; \xi_i \ge 0,
\end{equation}
has the dual
\begin{equation}
  \max_{\bm{\alpha}} \; \sum_{i=1}^{m} \alpha_i
    - \frac{1}{2} \sum_{i,j=1}^{m} \alpha_i \alpha_j \tilde{y}_i \tilde{y}_j\,
      \kappa(\bm{x}_i, \bm{x}_j)
  \quad \text{s.t.} \quad
  0 \le \alpha_i \le C, \;\; \textstyle\sum_i \alpha_i \tilde{y}_i = 0,
  \label{eq:svm-dual}
\end{equation}
with decision function
$f(\bm{x}) = \sum_i \alpha_i \tilde{y}_i \kappa(\bm{x}_i, \bm{x}) + b$
supported on the training points with $\alpha_i > 0$ (the \emph{support
vectors}). Two properties of \eqref{eq:svm-dual} matter here. First, the data
enter only through the \emph{Gram matrix}
$K_{ij} = \kappa(\bm{x}_i, \bm{x}_j)$ --- so any procedure that can estimate
pairwise kernel values, including a quantum computer estimating state
fidelities, can drive an otherwise entirely classical SVM. This is exactly the
quantum-kernel pipeline of Section~\ref{sec:qml-kernels}, and it is why the
SVM receives the most detailed treatment of any classical model in this
thesis. Second, the optimization is convex: given the Gram matrix, training
has no local-minimum or initialization pathologies, in sharp contrast to the
variational models of Section~\ref{sec:qml-vqa}.

The canonical general-purpose kernel is the radial basis function (RBF),
\begin{equation}
  \kappa_{\mathrm{RBF}}(\bm{x}, \bm{x}')
    = \exp\bigl( -\gamma \norm{\bm{x} - \bm{x}'}^2 \bigr),
  \label{eq:rbf}
\end{equation}
whose bandwidth $\gamma$ interpolates between nearly linear behavior
($\gamma \to 0$) and nearest-neighbor-like memorization ($\gamma \to \infty$).
The RBF-SVM is the single most important classical baseline in this thesis:
it is the like-for-like classical counterpart of the quantum-kernel SVM ---
same learning algorithm, different similarity function --- so any performance
difference is attributable to the kernel itself.

\section{Neural networks and gradient-based training}
\label{sec:ml-nn}

A multilayer perceptron (MLP) with one hidden layer of width $k$ computes
\begin{equation}
  f(\bm{x}) = \bm{w}_2^{\top}\, g\bigl( W_1 \bm{x} + \bm{b}_1 \bigr) + b_2,
\end{equation}
with an elementwise nonlinearity $g$ (ReLU in this thesis), giving
$dk + 2k + 1$ trainable parameters. Universal-approximation results guarantee
expressivity in the wide limit; what is actually achievable is decided by
optimization --- non-convex, performed by stochastic gradient descent on
mini-batches with gradients from backpropagation, in this work using the Adam
optimizer \citep{kingma2015adam} --- and by regularization ($\ell_2$ weight
penalties here) \citep{goodfellow2016deep}.

The MLP has a designated role in the fairness protocol: it is the
\emph{parameter-matched} classical mirror of the variational quantum
classifier. A VQC with $p$ trainable rotation angles is compared against MLPs
constrained to approximately $p$ parameters, so that RQ3's hybrid-vs-classical
comparisons oppose architectures of comparable capacity rather than comparable
marketing. Both model families also share an optimization story --- noisy
gradient estimates of a non-convex objective --- which makes their trainability
comparison (barren plateaus versus vanishing gradients,
Section~\ref{sec:qml-trainability}) meaningful.

\section{Ensembles: random forests and gradient boosting}
\label{sec:ml-ensembles}

Decision-tree ensembles are the strongest general-purpose baselines on small-%
to-medium tabular data --- precisely this thesis's data regime --- and their
absence is a recurring flaw in QML comparisons. A \emph{random forest}
\citep{breiman2001rf} averages trees grown on bootstrap resamples with random
feature subsetting at each split, reducing variance while keeping bias low.
\emph{Gradient boosting} \citep{chen2016xgboost} instead builds trees
sequentially, each fit to the gradient of the loss at the current ensemble's
predictions, with shrinkage and subsampling controlling capacity; this work
uses the XGBoost implementation. Neither model consumes the $[0,\pi]$ feature
scaling in any essential way --- trees are invariant to monotone feature maps
--- which makes them a useful control: a model family whose performance is
unaffected by the encoding-motivated preprocessing shared across the
comparison.

\section{Instance-based learning}
\label{sec:ml-knn}

$k$-nearest-neighbors classifies by majority vote among the $k$ closest
training points under a Minkowski metric. It is included as the minimal
``pure memorization'' baseline: no training, no parametric structure, purely
local generalization. Its behavior on the engineered synthetic tasks is
diagnostic --- it excels when local neighborhoods are label-pure (moons,
circles) and degrades when they are not (the sign-parity construction, where
distractor dimensions corrupt the metric) --- providing a reference point that
separates ``the task needs nonlinearity'' from ``the task needs
\emph{non-local} structure.''

\section{Evaluation methodology}
\label{sec:ml-evaluation}

\paragraph{Metrics.}
Accuracy is the headline metric on balanced tasks but degrades to
uninformativeness under imbalance. This thesis therefore reports four
complementary quantities per evaluation: \emph{accuracy};
\emph{balanced accuracy}, the mean of per-class recalls, robust to the
moderate imbalance of several included datasets; $F_1$, the harmonic mean of
precision and recall for the positive class; and \emph{ROC-AUC}, the
probability that a random positive instance outscores a random negative one,
which evaluates the score function independently of a threshold.

\paragraph{Resampling.}
A single train/test split yields a point estimate whose variance is unknown.
The protocol here uses $k$-fold \emph{stratified} cross-validation ($k = 5$),
which preserves class proportions per fold, repeated over five different
shuffling seeds --- 25 evaluations per model--dataset pair --- reporting means
and standard deviations across folds. Crucially, the split indices are
generated \emph{once}, stored, and reused by every model ever compared:
resampling variance then cancels in model comparisons, turning the folds into
matched pairs.

\paragraph{Hyperparameter optimization without leakage.}
Hyperparameters must be selected using training data only; selecting them on
test folds is the most common form of optimistic bias. The protocol fixes a
matched budget --- one 50-trial Optuna study \citep{akiba2019optuna} per
model--dataset pair, using the tree-structured Parzen estimator sampler with a
fixed seed, each trial scored by an internal 3-fold cross-validation confined
to the training portion of one designated outer fold --- after which the
selected configuration is frozen and evaluated on all 25 outer folds. The
residual weakness of this design (the tuning fold's training data reappears in
other folds' evaluations) is deliberate, tractability-motivated, and assessed
in Chapter~\ref{ch:discussion}; its key property is \emph{symmetry}: every
model, classical or quantum, is tuned by exactly the same procedure with
exactly the same budget \citep{bowles2024better}.

\paragraph{Statistical comparison.}
With 25 paired evaluations per model pair, per-dataset differences are tested
with the paired Wilcoxon signed-rank test, and comparisons across the whole
dataset suite follow the rank-based methodology of \citet{demsar2006comparisons},
including critical-difference diagrams for the headline comparison.
Significance levels are stated in advance in Chapter~\ref{ch:methodology};
non-significant differences are reported as such rather than narrativized.

\chapter{Quantum Computing Foundations}
\label{ch:quantum}

This chapter assembles the quantum-computing concepts on which
Chapter~\ref{ch:qml} builds: how quantum information is represented, how it is
transformed and measured, what today's hardware can realistically execute, and
where classical simulation of quantum circuits reaches its limits. The
treatment is deliberately focused; comprehensive accounts can be found in
\citet{nielsen2010quantum}.

\section{Qubits, superposition, and entanglement}
\label{sec:qc-qubits}

The elementary unit of quantum information is the \emph{qubit}, a two-level
quantum system. Its state is a unit vector in the two-dimensional complex
Hilbert space $\Hil_1 = \C^2$, written in Dirac notation as
\begin{equation}
  \ket{\psi} = \alpha \ket{0} + \beta \ket{1},
  \qquad \alpha, \beta \in \C, \quad
  \abs{\alpha}^2 + \abs{\beta}^2 = 1,
  \label{eq:qubit}
\end{equation}
where $\{\ket{0}, \ket{1}\}$ is the \emph{computational basis}. A state with
both $\alpha \neq 0$ and $\beta \neq 0$ is a \emph{superposition}: not a
probabilistic mixture of $\ket{0}$ and $\ket{1}$, but a distinct state whose
complex amplitudes can interfere. Global phase is unobservable, so a single
qubit state is fully described by two real angles,
\begin{equation}
  \ket{\psi} = \cos\tfrac{\vartheta}{2}\,\ket{0}
             + e^{i\varphi} \sin\tfrac{\vartheta}{2}\,\ket{1},
  \qquad \vartheta \in [0, \pi],\ \varphi \in [0, 2\pi),
\end{equation}
which identifies the state with a point on the \emph{Bloch sphere}. This
picture matters for machine learning because the most common way to load a
real-valued feature into a qubit --- \emph{angle encoding},
Section~\ref{sec:qml-encoding} --- is precisely to use the feature as a Bloch
rotation angle; the $[0,\pi]$ feature scaling adopted throughout this thesis
maps each feature onto a great arc of the sphere without wrapping.

A register of $n$ qubits lives in the tensor-product space
$\Hil_n = (\C^2)^{\otimes n} \cong \C^{2^n}$: a general state
\begin{equation}
  \ket{\psi} = \sum_{z \in \{0,1\}^n} c_z \ket{z},
  \qquad \sum_z \abs{c_z}^2 = 1
  \label{eq:nqubit}
\end{equation}
carries $2^n$ complex amplitudes. This exponential dimensionality is the
resource QML hopes to exploit --- and, simultaneously, the reason
$n \lesssim 25$-qubit systems remain classically simulable
(Section~\ref{sec:qc-simulability}).

A multi-qubit state that cannot be written as a tensor product of single-qubit
states is \emph{entangled}; the canonical example is the Bell state
$\ket{\Phi^+} = (\ket{00} + \ket{11})/\sqrt{2}$. Entanglement is generated by
two-qubit gates and is the structural ingredient that distinguishes genuinely
quantum feature maps and ansätze from products of independent single-qubit
operations. Whether that ingredient actually helps on standard learning tasks
is an empirical question this thesis probes directly with
entanglement-on/off ablations (Chapter~\ref{ch:experiments}), motivated by
benchmark findings that removing entanglement often costs little
\citep{bowles2024better}.

\section{Gates and circuits}
\label{sec:qc-gates}

Closed-system quantum evolution is unitary: a \emph{gate} on $k$ qubits is a
unitary $U \in \mathrm{U}(2^k)$ acting on the corresponding factors of
$\Hil_n$. The single-qubit gates used constantly in this thesis are the Pauli
operators
\begin{equation}
  X = \begin{pmatrix} 0 & 1 \\ 1 & 0 \end{pmatrix}, \quad
  Y = \begin{pmatrix} 0 & -i \\ i & 0 \end{pmatrix}, \quad
  Z = \begin{pmatrix} 1 & 0 \\ 0 & -1 \end{pmatrix},
\end{equation}
the Hadamard gate $H = \tfrac{1}{\sqrt{2}}\bigl(\begin{smallmatrix} 1 & 1 \\
1 & -1 \end{smallmatrix}\bigr)$, which maps computational-basis states to
uniform superpositions, and the parameterized rotations
\begin{equation}
  R_P(\theta) = e^{-i \theta P / 2}
              = \cos\tfrac{\theta}{2}\, I - i \sin\tfrac{\theta}{2}\, P,
  \qquad P \in \{X, Y, Z\}.
  \label{eq:rotation}
\end{equation}
Rotations are the workhorses of variational circuits: data values enter as
rotation angles (encoding), and trainable parameters $\bm{\theta}$ enter as
rotation angles (ansatz). The form \eqref{eq:rotation}, with generator
satisfying $P^2 = I$, is also what makes the parameter-shift rule of
Section~\ref{sec:qml-vqa} exact.

Entangling operations are typically the controlled-NOT
($\mathrm{CNOT}\,\ket{a,b} = \ket{a, a \oplus b}$) or controlled-$Z$ gates.
Together with single-qubit rotations they form a universal set: any unitary on
$n$ qubits can be approximated to arbitrary precision by a finite circuit over
this set \citep{nielsen2010quantum}.

A \emph{quantum circuit} is a sequence of gates applied to an initial state,
conventionally $\ket{0}^{\otimes n}$. Two cost measures matter throughout this
thesis. The \emph{gate count} determines simulation effort and hardware error
accumulation; the \emph{depth} --- the number of layers of gates that cannot be
parallelized --- determines how long a computation must survive decoherence.
On hardware, circuits must additionally be \emph{transpiled}: rewritten over
the device's native gate set and routed to respect its qubit-connectivity
graph, a step that can substantially inflate both counts and is logged
explicitly in the resource-accounting experiments (E9).

\section{Measurement and expectation values}
\label{sec:qc-measurement}

Information leaves a quantum computer only through measurement. Measuring all
qubits of state \eqref{eq:nqubit} in the computational basis returns bitstring
$z$ with probability $p(z) = \abs{c_z}^2$ (the \emph{Born rule}), and the state
collapses accordingly; a measurement is not a passive read-out but a
destructive sampling operation.

Machine-learning pipelines almost always consume \emph{expectation values} of
observables rather than raw bitstrings. For a Hermitian observable $M$ and
state $\ket{\psi}$, the quantity of interest is
\begin{equation}
  \langle M \rangle = \bra{\psi} M \ket{\psi},
\end{equation}
for example $M = Z_1$ (the Pauli-$Z$ on a designated readout qubit), whose
expectation in $[-1, 1]$ serves directly as a binary-classification score.
On real devices and shot-based simulators, $\langle M \rangle$ is estimated by
repeating the circuit $S$ times (``shots'') and averaging outcomes. The
estimator's standard error scales as
\begin{equation}
  \varepsilon \;\sim\; \sqrt{\frac{\Var[M]}{S}} \;=\; O\!\bigl(S^{-1/2}\bigr),
  \label{eq:shotnoise}
\end{equation}
so each additional decimal digit of precision costs a hundredfold more shots.
Equation~\eqref{eq:shotnoise} is the root of two phenomena central to this
thesis: the \emph{shot-noise regime} of experiment E8, where model quality
degrades gracefully as $S$ shrinks, and the interaction of sampling error with
the \emph{concentration} effects of Section~\ref{sec:qml-trainability}, where
signals that decay exponentially in the qubit count fall below any affordable
shot budget.

\section{The NISQ era: noise and error mitigation}
\label{sec:qc-nisq}

Current devices --- tens to a few hundred qubits, no fault tolerance ---
define the \emph{noisy intermediate-scale quantum} (NISQ) era
\citep{preskill2018nisq}. Physical qubits decohere: excited states relax
toward $\ket{0}$ on a timescale $T_1$, and phase coherence between $\ket{0}$
and $\ket{1}$ decays on a timescale $T_2$. Gates are imperfect, with two-qubit
gates typically an order of magnitude noisier than single-qubit ones, and
readout itself misassigns outcomes with percent-level probability.

Noise is modeled by quantum channels acting on density matrices
$\rho = \ketbra{\psi}{\psi}$ (and their mixtures). The depolarizing channel
with error probability $p$,
\begin{equation}
  \mathcal{E}_{\mathrm{dep}}(\rho)
    = (1 - p)\,\rho + \frac{p}{3}\left( X\rho X + Y\rho Y + Z\rho Z \right),
\end{equation}
is the generic ``white noise'' model; bit-flip, phase-flip, and
amplitude-damping channels capture more specific error mechanisms, and
calibrated device models compose gate-level channels with measured $T_1$,
$T_2$, gate-error, and readout-error parameters of a physical backend. The
noise ladder of experiment E8 uses exactly this hierarchy: exact statevector
$\rightarrow$ finite shots $\rightarrow$ calibrated device-noise simulation
$\rightarrow$ real hardware.

Full quantum error \emph{correction} is beyond NISQ budgets, but lightweight
error \emph{mitigation} is not. This thesis uses measurement-error (readout)
mitigation --- characterizing the readout-confusion matrix and applying its
(regularized) inverse to measured distributions --- for the hardware runs, and
notes zero-noise extrapolation as a standard alternative; heavier schemes are
out of scope. Mitigation reduces bias in expectation values at the price of
increased variance, i.e., more shots \citep{preskill2018nisq}.

\section{Limits of classical simulability}
\label{sec:qc-simulability}

Storing the state vector \eqref{eq:nqubit} of $n$ qubits requires $2^n$
complex amplitudes: at double precision, roughly $16 \times 2^n$ bytes ---
about 4\,GB at $n = 28$ and about 1\,TB at $n = 36$. Statevector simulation of
the circuit sizes in this thesis ($n \le 20$) is therefore comfortable on a
laptop, which is precisely why the experimental programme can run locally,
exactly, and reproducibly. Restricted circuit families admit far larger
simulations --- Clifford circuits via the stabilizer formalism, and
low-entanglement circuits via tensor-network contraction
\citep{nielsen2010quantum} --- and modern classical surrogate and
dequantization results (Section~\ref{sec:related-surrogates}) push the
simulable frontier further for many structured QML models specifically.

Two consequences frame everything that follows. First, the honest-framing
position of Section~\ref{sec:intro-framing} is not a caveat but a mathematical
fact: at the scales studied here, quantum models \emph{are} classical
computations, executed through an unusual parameterization. Second,
simulability is an asset for methodology: exact simulation provides
noise-free ground truth against which every shot, noise, and hardware effect
in E8 can be measured --- a separation of concerns impossible in
hardware-first studies.

\chapter{Quantum Machine Learning}
\label{ch:qml}

This chapter develops the quantum and hybrid model families that
Chapters~\ref{ch:methodology} and \ref{ch:experiments} put to the test. The
running theme is that every design choice --- how data enter the circuit, how
the trainable part is shaped, what is measured, and how parameters are
optimized --- has consequences that are now reasonably well understood
theoretically, and that these consequences (expressivity spectra, kernel
equivalences, barren plateaus, concentration) directly generate the hypotheses
of the experimental programme.

\section{A taxonomy of quantum machine learning}
\label{sec:qml-taxonomy}

Following the scheme popularized by \citet{schuld2021book}, QML approaches are
classified by whether the \emph{data} and the \emph{processing device} are
classical (C) or quantum (Q). \textbf{CC} is classical ML, possibly
quantum-inspired. \textbf{CQ} --- classical data, quantum processing --- is
the subject of this thesis: feature vectors are encoded into quantum states,
processed by circuits, and read out by measurement. \textbf{QC} uses
classical ML to support quantum experiments (calibration, tomography,
control), and \textbf{QQ} learns on inherently quantum data, the setting where
the strongest advantage arguments live \citep{huang2022advantage}; both are
outside the experimental scope and appear only in the outlook.

Within CQ, this thesis studies the two dominant paradigms --- \emph{quantum
kernel methods} (Section~\ref{sec:qml-kernels}) and \emph{variational quantum
circuits} (Sections~\ref{sec:qml-vqa}--\ref{sec:qml-architectures}) --- plus
\emph{hybrid architectures} that compose them with classical networks
(Section~\ref{sec:qml-hybrid}). The taxonomy earns its place beyond
bookkeeping with one observation: a quantum kernel inside a classical SVM is
already a hybrid CQ pipeline, so ``quantum versus hybrid versus classical'' is
a spectrum of composition choices, not three disjoint species.

\section{Data encoding}
\label{sec:qml-encoding}

A CQ model must first map $\bm{x} \in \R^d$ to an $n$-qubit state
$\ket{\phi(\bm{x})} = \enc{\bm{x}} \ket{0}^{\otimes n}$ via an encoding
(feature-map) circuit $\enc{\bm{x}}$. This choice is not preprocessing
trivia: it fixes the geometry of the data in Hilbert space, hence the kernel
the model can realize and the function class it can express
\citep{schuld2019hilbert, schuld2021fourier, larose2020encodings}.

\paragraph{Basis encoding} writes a binary string into the computational
basis, $\bm{b} \mapsto \ket{\bm{b}}$. It is qubit-hungry for continuous data
and is used here only conceptually.

\paragraph{Angle encoding} uses one feature per qubit as a rotation angle,
\begin{equation}
  \enc{\bm{x}} = \bigotimes_{j=1}^{d} R_{Y}(x_j),
  \qquad
  \ket{\phi(\bm{x})} = \bigotimes_{j=1}^{d}
    \Bigl( \cos\tfrac{x_j}{2}\ket{0} + \sin\tfrac{x_j}{2}\ket{1} \Bigr),
  \label{eq:angle-encoding}
\end{equation}
requiring $n = d$ qubits and depth 1. Because the rotation is $2\pi$-periodic,
features must be confined to a fixed interval to avoid wrap-around aliasing;
this motivates the thesis-wide convention of min--max scaling every feature to
$[0, \pi]$, fit on training folds only (Chapter~\ref{ch:methodology}).
Unentangled angle encoding produces a product state, and its induced kernel
factorizes into a product of per-feature $\cos$ terms --- a useful reminder
that ``quantum'' does not yet mean ``hard to simulate.''

\paragraph{Entangled feature maps.} Expressive power beyond products comes
from interleaving encodings with entanglers. The canonical example is the
ZZ (IQP-style) feature map of \citet{havlicek2019kernels},
\begin{equation}
  \enc{\bm{x}}
  = \Bigl[ U_{\Phi}(\bm{x})\, H^{\otimes n} \Bigr]^{L},
  \qquad
  U_{\Phi}(\bm{x}) = \exp\Bigl( i \sum_{j} \phi_j(\bm{x})\, Z_j
      + i \sum_{j < k} \phi_{jk}(\bm{x})\, Z_j Z_k \Bigr),
  \label{eq:zz-map}
\end{equation}
with the standard choices $\phi_j(\bm{x}) = x_j$ and
$\phi_{jk}(\bm{x}) = (\pi - x_j)(\pi - x_k)$, repeated $L$ times. The
product terms couple pairs of features through entangling phases ---
conjectured hard to simulate classically at depth $L \ge 2$ --- and make
parity-like feature interactions natively available, which is precisely why
the sign-parity benchmark task of Chapter~\ref{ch:methodology} is the
designated stress test for this family.

\paragraph{Amplitude encoding} stores $\bm{x}$ (padded, normalized) in the
amplitudes of \eqref{eq:nqubit}: $2^n$ features fit in $n$ qubits, but state
preparation generally costs depth exponential in $n$ and the encoding is
nonlinear in $\bm{x}$ through normalization. It appears in this thesis only
as a comparison point in the encoding ablation (E4).

\paragraph{Data re-uploading} repeats the encoding several times, interleaved
with trainable blocks \citep{perezsalinas2020reuploading}:
\begin{equation}
  \ket{\phi_{\bm{\theta}}(\bm{x})}
  = \prod_{l=1}^{L} \Bigl[ W_l(\bm{\theta}_l)\, \enc{\bm{x}} \Bigr]
    \ket{0}^{\otimes n}.
  \label{eq:reuploading}
\end{equation}
Re-uploading is not a redundancy: it provably enlarges the model's frequency
content, as the following Fourier picture makes precise.

\paragraph{The Fourier lens on expressivity.}
\citet{schuld2021fourier} showed that any model of the form
$f_{\bm{\theta}}(\bm{x}) = \bra{\phi_{\bm{\theta}}(\bm{x})} M
\ket{\phi_{\bm{\theta}}(\bm{x})}$ whose encoding gates are generated by fixed
Hamiltonians, $e^{-i x_j H_j}$, is a generalized trigonometric polynomial
\begin{equation}
  f_{\bm{\theta}}(\bm{x})
  = \sum_{\bm{\omega} \in \Omega} c_{\bm{\omega}}(\bm{\theta}, M)\,
    e^{i \bm{\omega}^{\top} \bm{x}},
  \label{eq:fourier}
\end{equation}
whose frequency set $\Omega$ is determined \emph{entirely by the encoding} ---
the pairwise differences of eigenvalue sums of the $H_j$ --- while the
trainable parts and the observable only shape the coefficients
$c_{\bm{\omega}}$. For Pauli-rotation encodings, one upload per feature yields
frequencies $\{-1, 0, 1\}$ per dimension; $L$ uploads extend this to
$\{-L, \dots, L\}$. Three consequences organize much of what follows:
(i) a single-upload angle-encoded model is, in each feature, a degree-one
trigonometric function --- expressivity claims must be checked against this
sobering baseline; (ii) re-uploading, not deeper ansätze, is the lever that
adds frequencies; and (iii) expressivity is decoupled from trainability ---
enlarging $\Omega$ says nothing about whether $c_{\bm{\omega}}$ can be
efficiently fit (Section~\ref{sec:qml-trainability}). Experiment E4
demonstrates \eqref{eq:fourier} directly by fitting one-dimensional target
functions with increasing upload counts.

\section{Variational quantum algorithms}
\label{sec:qml-vqa}

The variational paradigm \citep{cerezo2021vqa} builds models from three
components: an encoding $\enc{\bm{x}}$, a trainable ansatz $\ansatz$, and a
measured observable $M$,
\begin{equation}
  f_{\bm{\theta}}(\bm{x})
  = \bra{0^{\otimes n}} \enc{\bm{x}}^{\dagger}\, W^{\dagger}(\bm{\theta})\,
    M\, W(\bm{\theta})\, \enc{\bm{x}} \ket{0^{\otimes n}},
  \label{eq:vqc-model}
\end{equation}
trained by a classical optimizer minimizing an empirical loss
$C(\bm{\theta}) = \tfrac{1}{m} \sum_i \ell\bigl( f_{\bm{\theta}}(\bm{x}_i),
y_i \bigr)$ --- a quantum--classical loop in which the quantum device only
evaluates \eqref{eq:vqc-model}. In this thesis $M$ is a single-qubit
$Z$ observable unless stated otherwise, making
$f_{\bm{\theta}} \in [-1, 1]$ a signed classification score
\citep{mitarai2018qcl, farhi2018qnn, schuld2020circuitcentric}.

\paragraph{Ansatz families.} The experiments use two representative ansätze:
\emph{hardware-efficient} layers of single-qubit rotations plus a fixed
entangling pattern \citep{kandala2017hardware}, chosen for shallow depth on
real devices; and \emph{strongly entangling layers}
\citep{schuld2020circuitcentric}, with all-to-all-leaning connectivity and
three rotation angles per qubit per layer. Depth is swept explicitly (E3),
since it controls both capacity and every pathology of
Section~\ref{sec:qml-trainability}.

\paragraph{Exact gradients: the parameter-shift rule.}
For gates of the form \eqref{eq:rotation} --- generators with two eigenvalues
$\pm\tfrac{1}{2}$ --- the derivative of \eqref{eq:vqc-model} with respect to
the $k$-th angle is \emph{exactly}
\begin{equation}
  \frac{\partial f_{\bm{\theta}}}{\partial \theta_k}
  = \frac{1}{2} \Bigl[
      f_{\bm{\theta} + \frac{\pi}{2} \bm{e}_k}
    - f_{\bm{\theta} - \frac{\pi}{2} \bm{e}_k}
    \Bigr],
  \label{eq:parameter-shift}
\end{equation}
two evaluations of the same circuit at shifted parameters
\citep{mitarai2018qcl, schuld2021fourier} --- not a finite-difference
approximation, and compatible with hardware since only expectation values are
needed. The cost, however, scales linearly in the parameter count \emph{per
gradient step} (two circuit evaluations per parameter, each itself averaged
over shots), which is why gradient-based VQC training on hardware is
prohibitive and why this thesis trains exclusively in simulation, reserving
hardware for inference (Chapter~\ref{ch:methodology}).

\paragraph{Optimizers.} Experiment E3 compares Adam \citep{kingma2015adam} on
parameter-shift gradients with two gradient-light alternatives standard in
the VQA literature: SPSA \citep{spall1992spsa}, which estimates a descent
direction from two evaluations regardless of dimension and is naturally robust
to shot noise, and the derivative-free trust-region method COBYLA
\citep{powell1994cobyla}. The optimizer axis is treated as a first-class
design factor (RQ2), not a nuisance setting.

\section{Quantum kernel methods}
\label{sec:qml-kernels}

The encoding circuit alone already defines a learning method. The
\emph{fidelity quantum kernel} associated with $\enc{\cdot}$ is
\begin{equation}
  \qkernel(\bm{x}, \bm{x}')
  = \abs{ \braket{\phi(\bm{x})}{\phi(\bm{x}')} }^2
  = \abs{ \bra{0^{\otimes n}} \enc{\bm{x}}^{\dagger} \enc{\bm{x}'}
          \ket{0^{\otimes n}} }^2,
  \label{eq:fidelity-kernel}
\end{equation}
a positive-semidefinite similarity measure computable on a quantum device by
preparing $\enc{\bm{x}}^{\dagger} \enc{\bm{x}'} \ket{0^{\otimes n}}$ and
measuring the probability of returning to $\ket{0^{\otimes n}}$ (the
``inversion test''), at circuit depth twice that of the encoding
\citep{havlicek2019kernels, schuld2019hilbert}. The QSVM pipeline estimates
the Gram matrix $K_{ij} = \qkernel(\bm{x}_i, \bm{x}_j)$ --- $O(m^2)$ circuit
families, each averaged over shots --- and hands it to the convex classical
SVM dual \eqref{eq:svm-dual}. Against the RBF-SVM baseline
\eqref{eq:rbf}, the QSVM differs in exactly one component, which is what
makes the E2 comparison clean.

\paragraph{Quantum models are kernel methods.}
\citet{schuld2021kernels} sharpened the connection: the class of models
\eqref{eq:vqc-model}, viewed as functions of the encoded density matrix
$\rho(\bm{x}) = \ketbra{\phi(\bm{x})}{\phi(\bm{x})}$, is \emph{linear} in
$\rho(\bm{x})$, i.e., a linear model in the feature space of density matrices
with inner product $\Tr[\rho(\bm{x}) \rho(\bm{x}')] = \qkernel(\bm{x},
\bm{x}')$. By representer-theorem reasoning, the loss-optimal model in that
class is a kernel expansion over the training data --- so a trained VQC can,
at best, match the regularized kernel machine built from
\eqref{eq:fidelity-kernel}, and variational training is better understood as
implicit (possibly advantageous, possibly lossy) regularization within the
kernel's hypothesis space. This result frames E2 and E3 as two views of one
object and cautions against narratives in which the ansatz ``adds power''
beyond the encoding.

\paragraph{Kernel diagnostics.} Two quantities computed in E2 connect kernels
to outcomes before any SVM is trained: \emph{kernel--target alignment}, the
normalized inner product between the Gram matrix and the ideal label kernel
$\bm{\tilde{y}} \bm{\tilde{y}}^{\top}$, a cheap predictor of learnability; and
the \emph{spectrum} of $K$, whose flatness signals the concentration pathology
below. The \emph{projected quantum kernel} of \citet{huang2021power} ---
built from reduced single-qubit density matrices,
$\kappa^{\mathrm{PQ}}(\bm{x}, \bm{x}') = \exp\bigl( -\gamma \sum_k
\norm{\rho_k(\bm{x}) - \rho_k(\bm{x}')}_F^2 \bigr)$ --- is included as the
literature's main remedy: by projecting before comparing, it tempers the
exponential-dimension geometry that drives both concentration and poor
generalization \citep{kubler2021inductive}.

\section{Model architectures}
\label{sec:qml-architectures}

Three circuit architectures span the variational side of the comparison.

\paragraph{VQC.} The plain variational quantum classifier
\eqref{eq:vqc-model}, instantiated over the grid \{angle, ZZ\} encodings
$\times$ \{hardware-efficient, strongly entangling\} ansätze $\times$ depth
$\{1, 3, 5\}$ --- the E3 design. Parameter counts ($3nd_{\mathrm{ansatz}}$ for
strongly entangling layers on $n$ qubits) are logged for every run and define
the MLP matching budgets.

\paragraph{Data re-uploading classifier.} The single- or few-qubit model
\eqref{eq:reuploading} of \citet{perezsalinas2020reuploading}, universal for
classification on one qubit in the many-upload limit and, per
\eqref{eq:fourier}, the cleanest experimental handle on the
frequency-enrichment story (E4). It is also among the strongest small quantum
models in independent benchmarks \citep{bowles2024better}, at minimal
simulation cost.

\paragraph{QCNN.} The quantum convolutional network of \citet{cong2019qcnn}
alternates translationally-shared two-qubit ``convolution'' unitaries with
measurement-free ``pooling'' that discards half the qubits per layer, giving
$O(\log n)$ depth and parameter counts logarithmic in the input size; the
classical-data variant follows \citet{hur2022qcnn}. QCNNs are the image-tier
quantum contender (8 qubits over PCA-compressed pixel features) and are of
special theoretical interest because the architecture is among the few with
evidence of \emph{avoiding} barren plateaus \citep{larocca2025barren} --- at
the price of falling within classically simulable families
\citep{cerezo2023simulability}, a tension E5/E9 quantify in practice.

\section{Trainability: barren plateaus and concentration}
\label{sec:qml-trainability}

The central obstruction to scaling variational models is not expressivity but
\emph{trainability}. \citet{mcclean2018barren} showed that for ansätze random
enough to approximate unitary 2-designs, gradients of cost functions like
\eqref{eq:vqc-model} vanish in probability:
\begin{equation}
  \E\bigl[ \partial_{\theta_k} C \bigr] = 0,
  \qquad
  \Var\bigl[ \partial_{\theta_k} C \bigr] \in O\!\bigl( b^{-n} \bigr)
  \text{ for some } b > 1,
  \label{eq:barren}
\end{equation}
the \emph{barren plateau}: with overwhelming probability over random
initializations, the loss landscape is exponentially flat in the qubit count,
and estimating a useful gradient direction through shot noise
\eqref{eq:shotnoise} costs exponentially many measurements.
\citet{cerezo2021costfunction} refined the picture: with shallow
(logarithmic-depth) ansätze, \emph{local} observables (as used throughout this
thesis) admit gradients decaying at worst polynomially, while \emph{global}
observables plateau at any depth --- making cost design a first-order modeling
decision. The modern synthesis \citep{larocca2025barren} identifies
randomness, entanglement growth, noise, and cost globality as interchangeable
sources of the same exponential concentration; strategies that provably escape
it (structured ansätze such as QCNNs, clever initializations, problem-tailored
symmetry) tend to land in circuit families that are classically simulable
\citep{cerezo2023simulability} --- arguably the sharpest form of the
quantum-utility dilemma. Experiment E6 reproduces \eqref{eq:barren} directly,
measuring gradient variance versus qubit count ($2$--$16$) and depth, for
local versus global costs.

Quantum kernels suffer a mirror-image pathology: \emph{exponential
concentration} of kernel values \citep{thanasilp2024concentration}. For
sufficiently expressive (or deep, or noisy) encodings, $\qkernel(\bm{x},
\bm{x}')$ concentrates around a data-independent constant with deviations
exponentially small in $n$ --- so the Gram matrix approaches a scaled
identity, generalization collapses, and resolving the residual signal requires
exponentially many shots. Bandwidth-style fixes (rescaling feature ranges) and
projected kernels trade expressivity back for signal
\citep{kubler2021inductive, huang2021power}. E2 measures kernel-value spread
versus qubit count on this thesis's datasets. Finally, on the generalization
side, uniform bounds scaling with trainable-gate counts \citep{caro2022generalization}
--- and their empirical refutation as a complete story
\citep{gilfuster2024rethinking} --- plus the effective-dimension perspective
\citep{abbas2021power} justify treating capacity claims empirically (RQ5)
rather than as settled theory.

\section{Hybrid architectures}
\label{sec:qml-hybrid}

Hybrid models compose classical and quantum components under end-to-end or
staged training; automatic differentiation through quantum nodes
\citep{bergholm2018pennylane} makes the composition practical. Three designs
enter the comparison (E5), each with a matched classical control.

\paragraph{Compression front-ends.} PCA (or a small autoencoder) reduces
$d$-dimensional inputs to the qubit budget before a VQC. This is the standard
route for image data and doubles as the fair way to give quantum models
high-dimensional inputs at all; the control is the same compressed features
fed to classical heads.

\paragraph{Classical feature extractor with quantum head.} A small CNN
produces $k$ features consumed by a $k$-qubit VQC head, trained jointly; the
control replaces the quantum head with a dense head parameter-matched to the
VQC's angle count. This is the cleanest hybrid-versus-classical contrast in
the thesis: identical data, identical backbone, matched capacity, single
varying component.

\paragraph{Quantum transfer learning.} Following \citet{mari2020transfer}, a
frozen pretrained backbone feeds a trainable ``dressed quantum circuit'' ---
classical linear map in, PQC, classical linear map out. Notably, this family
ranked among the stronger quantum-adjacent models in independent benchmarking
\citep{bowles2024better}, with the open question --- how much of its
performance is the dressing rather than the quantum core --- addressed by the
parameter-matched dressed-\emph{classical} control.

Throughout, the taxonomy of Section~\ref{sec:qml-taxonomy} keeps the
bookkeeping honest: the quantum kernel SVM is a CQ hybrid with quantum
similarity and classical optimization; the dressed circuit is a CQ hybrid with
classical representation learning and quantum decision surface. What the
experiments compare is therefore not ``quantum versus classical'' as slogans,
but specific placements of a quantum component inside otherwise-matched
pipelines --- exactly the resolution at which RQ3 is answerable.

\chapter{Related Work: Comparative Studies and the Advantage Debate}
\label{ch:related}

The empirical part of this thesis enters a conversation with four strands of
literature: systematic benchmarks of quantum models, the ``power of data''
argument, classical surrogates and dequantization results, and the small
island of provable separations. This chapter reviews each and closes by
positioning the present work.

\section{Systematic benchmark studies}
\label{sec:related-benchmarks}

For most of QML's short history, empirical claims rested on one-dataset,
one-baseline demonstrations. \citet{bowles2024better} changed the standard:
benchmarking twelve popular quantum models (and classical controls) across
structured dataset families with systematic hyperparameter search, they found
--- almost uniformly --- that well-tuned classical models match or beat the
quantum contenders on standard classical-data tasks, that out-of-the-box
quantum models are highly sensitive to hyperparameters, and, strikingly, that
\emph{removing entanglement} from several quantum architectures often does not
degrade accuracy, suggesting the quantum ingredient frequently is not what
carries the performance. Their methodological catalogue of the field's
recurring sins --- tuned-quantum versus default-classical comparisons, absent
strong baselines (tree ensembles in particular), inconsistent splits, and
regime mixing --- reads as a checklist of what this thesis's fairness protocol
(Chapter~\ref{ch:methodology}) is designed to preclude. In the kernel corner
specifically, \citet{schnabel2025scrutiny} subjected fidelity- and
projected-kernel pipelines to the same discipline and reached similarly
sobering conclusions, while emphasizing bandwidth and preprocessing as
first-order factors --- directly informing the E2 design here. The
replication-style component of this thesis (entanglement on/off ablations, the
default-versus-tuned deltas of E1) is a deliberate test of whether these
findings transfer to a different dataset suite and protocol.

\section{The power of data}
\label{sec:related-powerofdata}

\citet{huang2021power} reframed the advantage question: many quantum
computations become classically learnable once a classical model has access
to \emph{data} generated by the process --- training data can substitute for
quantum computation. They provide a geometric, dataset-specific test (based on
the model-independent \emph{geometric difference} between kernel matrices)
that upper-bounds how much any quantum kernel can outperform a classical one
on a given dataset, and introduce the projected quantum kernel to engineer
tasks where the gap is provably large. Two lessons shape this thesis. First,
dataset-conditional diagnostics (alignment, geometric difference) belong in
the pipeline \emph{before} performance claims --- E2 computes them. Second,
the datasets on which quantum kernels could plausibly excel are special by
construction; finding an advantage on generic UCI-style data would contradict
theory, so the honest expectation for RQ1 is parity at best, and the
scientifically interesting content lies in \emph{where and why} deviations
occur (the sign-parity task being this thesis's engineered candidate).

\section{Classical surrogates and dequantization}
\label{sec:related-surrogates}

A complementary deflationary line shows that trained quantum models often
admit efficient \emph{classical surrogates}. \citet{schreiber2023surrogates}
exploit exactly the Fourier representation \eqref{eq:fourier}: since a
re-uploading model is a trigonometric polynomial with an
encoding-determined frequency set, one can fit that function class directly,
classically, from evaluations --- matching the quantum model's behavior
without a quantum device. More broadly, \citet{cerezo2023simulability} argue
that the very structures that provably avoid barren plateaus (shallow local
circuits, QCNNs, symmetry-restricted ansätze) confine the dynamics to
polynomially large subspaces that are classically simulable --- crystallizing
a dilemma: at NISQ scale, \emph{trainable} tends to imply \emph{simulable}.
For a thesis running at $\le 20$ qubits this is not a threat but a framing
device (Section~\ref{sec:intro-framing}): the object of study is the behavior
of these models, with simulability guaranteed either way.

\section{Provable separations and quantum data}
\label{sec:related-separations}

The existence proofs live elsewhere. \citet{liu2021rigorous} construct a
classification task from the discrete-logarithm problem for which a fidelity
quantum kernel learns efficiently while any efficient classical learner would
break a standard cryptographic assumption --- a rigorous, end-to-end
separation, on a deliberately cryptographic (not naturally occurring) task;
it defines the stretch experiment E10. On genuinely \emph{quantum} data, the
case is stronger still: \citet{huang2022advantage} demonstrate, on real
hardware, exponential advantages in learning from quantum-enhanced
experiments (QQ in the taxonomy). Both results delimit rather than contradict
this thesis: they identify where advantage is theoretically at home ---
structured algebraic tasks and quantum data --- and thereby explain why a
matched-conditions comparison on classical benchmarks (this work) should be
read as mapping the terrain of \emph{competitiveness}, not searching for
advantage where theory says none should live.

\section{Positioning of this thesis}
\label{sec:related-positioning}

Relative to this literature, the present work makes four commitments that,
combined, are not standard even in the strongest prior benchmarks:
(i)~\emph{one} stored split system --- 5-fold stratified $\times$ five seeds,
generated once, shared by every model of every family, forever; (ii)~a
\emph{matched and audited} tuning budget (one 50-trial study per
model--dataset pair under an identical procedure, with the full trial log
published); (iii)~\emph{parameter matching} for every NN-style comparison and
\emph{regime separation} (exact/shots/noisy/hardware never pooled) for every
reported number; and (iv)~a dataset suite deliberately spanning saturated to
hard, including an engineered parity-structured task that instantiates the
theory-motivated best case for entangling feature maps. The intended
contribution is thus less a verdict than a reproducible measurement
instrument --- in the spirit of \citet{bowles2024better}, extended with the
budget-matching and regime discipline their audit calls for.

\chapter{Methodology}
\label{ch:methodology}

\todonote{Full prose deferred until the experimental programme completes
(user decision 2026-07-07: theory first, experiments last). The protocol
itself is already frozen and implemented in the repository; this skeleton
records the chapter plan and the facts that will not change.}

\section{Datasets and preprocessing}
\label{sec:method-datasets}
% - 10-dataset suite (synthetic / tabular / image tiers) with the rationale per
%   tier; table generated from data/ registry metadata (no hand-typed numbers).
% - Difficulty spectrum by design: saturated (banknote) -> hard (pima,
%   xor_gauss); xor_gauss construction spelled out (4-D sign parity + 2
%   distractor dims) as the engineered best case for entangling feature maps.
% - Shared preprocessing: per-feature min-max to [0, pi] fit on training folds
%   only; identical features for all model families; binarization choices
%   (iris versicolor/virginica; wine hardest pair by committed probe; digits
%   3-vs-5) with their audit trail.

\section{The fairness protocol}
\label{sec:method-protocol}
% - Stored splits: 5-fold stratified x seeds 0-4, generated once, checksummed,
%   shared by every model; runner refuses to re-split.
% - Matched tuning budget: one 50-trial TPE(seed 0) study per (model, dataset);
%   trials scored by 3-fold stratified CV on the training portion of s0_f0;
%   locked config evaluated on all 25 folds; NON-NESTED by design -> residual
%   overlap stated as limitation here and quantified in ch08.
% - Parameter matching for NN-style comparisons; n_params logged per run.
% - Regime separation: exact / shots / noisy / hardware / classical never mixed.

\section{Model zoo and hyperparameter spaces}
\label{sec:method-models}
% - Classical: logreg, SVM-linear, SVM-RBF, kNN, RF, XGBoost, small MLP
%   (search spaces <= 6 hyperparameters each; table auto-generated from
%   models/classical.py).
% - Quantum: fidelity-kernel QSVM (angle, ZZ); VQC grid; re-uploading; QCNN.
% - Hybrid: PCA->VQC; CNN->quantum head vs matched dense head; dressed circuit.

\section{Execution regimes, simulators, and hardware}
\label{sec:method-regimes}
% - lightning.qubit statevector (exact); Aer shot-based; calibrated fake
%   backends / NoiseModel.from_backend; IBM Open Plan validation runs (E8),
%   inference-only, with the QPU budget tactics from groundwork §10.

\section{Metrics, statistics, and reproducibility}
\label{sec:method-stats}
% - accuracy, balanced accuracy, F1, ROC-AUC per (seed, fold); mean ± std over
%   the 25 shared evaluations; paired Wilcoxon + critical-difference diagrams
%   (Demsar) with alpha stated up front.
% - Append-only results CSVs with config hashes, seeds, split ids, regime,
%   n_params, timings, git commit; every figure/table regenerated by a script
%   in analysis/.

\chapter{Experiments and Results}
\label{ch:experiments}

\todonote{Written last (user decision 2026-07-07). One section per experiment
family, each in the fixed format: hypothesis (pre-registered in
docs/groundwork.md §9) $\to$ setup $\to$ results $\to$ discussion. Every
figure/table is included from analysis/output/, never typed.}

\section{E1: Classical baseline sweep}
\label{sec:exp-e1}
% COMPLETE in the repository (results/e1_baselines*.csv; tables in
% analysis/output/). Prose pending. Facts to carry into the text:
% - 7 models x 10 datasets x 50-trial matched tuning; defaults variant too.
% - \resultref{master table: analysis/output/e1_master_tuned_balanced_accuracy}
% - Sanity signatures: linear models at chance on circles/xor_gauss; tuning
%   gains small except mlp/circles (+0.23) and svm_rbf/xor_gauss (+0.137).
% - Difficulty spectrum: banknote saturated (1.000) ... pima 0.727,
%   xor_gauss 0.690 -> headroom argument for E2+.

\section{E2: Quantum kernels}
\label{sec:exp-e2}
% QSVM (angle, ZZ, depth 1-2) vs RBF-SVM on identical splits; Gram matrices;
% kernel-target alignment; geometric difference (Huang et al.); concentration
% vs qubit count. Named hypothesis: ZZ map on xor_gauss.

\section{E3: VQC grid}
\label{sec:exp-e3}
% {angle, ZZ} x {hardware-efficient, strongly entangling} x depth {1,3,5};
% Adam(parameter-shift) vs SPSA vs COBYLA; training curves; MLP parameter
% matching enforced here.

\section{E4: Encoding ablation and Fourier expressivity}
\label{sec:exp-e4}
% Same ansatz, different encodings; 1-D function fit demo showing accessible
% frequency spectrum growth with re-uploads (Schuld et al. 2021).

\section{E5: Hybrid architectures}
\label{sec:exp-e5}
% PCA->VQC vs raw VQC; CNN->quantum head vs matched dense head; dressed
% quantum circuit transfer learning; controls parameter-matched.

\section{E6: Trainability study}
\label{sec:exp-e6}
% Gradient variance vs qubits (2-16) and depth; local vs global cost;
% log-scale variance plot reproducing barren-plateau scaling.

\section{E7: Sample efficiency}
\label{sec:exp-e7}
% Learning curves (n_train in {20..1000}) for best classical / quantum /
% hybrid per dataset, on training-side subsamples of the stored splits.

\section{E8: Noise ladder and hardware validation}
\label{sec:exp-e8}
% exact -> shots (1k-10k) -> calibrated device noise -> IBM QPU (2-3 best
% small models, readout mitigation); degradation waterfall; backend,
% calibration date, and job ids logged.

\section{E9: Resource accounting}
\label{sec:exp-e9}
% Wall-clock, parameter counts, circuit depth, total shots; cost-vs-accuracy
% scatter across all families.

\chapter{Discussion}
\label{ch:discussion}

\todonote{Written after ch07. Planned structure below.}

\section{Synthesis across research questions}
% RQ1-RQ5 answered explicitly, each pointing at the figures/tables that carry
% the evidence; where results are null, say so plainly.

\section{Limitations and threats to validity}
% - Scale: <= 20 qubits, everything classically simulable (by design).
% - Non-nested tuning: residual overlap between the tuning fold's training
%   data and other folds' evaluations; direction and rough size of the bias.
% - Budget sensitivity: 50-trial ceiling may favor models with smaller search
%   spaces; mitigated by the "size philosophy" cap, not eliminated.
% - Dataset scope: standard benchmarks + engineered synthetics; no natural
%   quantum data.
% - Hardware sample: few QPU runs, one vendor, one calibration window.

\section{When are quantum or hybrid approaches worth considering?}
% Decision-oriented reading of the evidence: task structure (parity-like
% interactions), data regime (small n), cost tolerance, and the honest
% competitiveness picture at NISQ scale.

\chapter{Conclusion and Future Work}
\label{ch:conclusion}

\todonote{Written last. Planned content: one-paragraph restatement of the
question and method; the headline empirical findings per RQ; the
methodological contribution (fairness protocol + reproducible artifact) as a
result in its own right; future work: multi-class extensions of the best
three models, Tier-3 quantum-native datasets, the discrete-log separation
task (E10) if not run, larger hardware validation under error mitigation, and
tracking the trainable-implies-simulable frontier as theory develops.}


\appendix
\chapter{Circuit Diagrams}
\label{app:circuits}
\todonote{Quantikz diagrams for: angle encoding, ZZ feature map (depth 1-2),
hardware-efficient and strongly entangling ansatz layers, re-uploading
classifier, QCNN convolution/pooling blocks, the inversion test for fidelity
kernels. Added when E2-E5 fix the exact circuits.}

\chapter{Full Hyperparameter Tables}
\label{app:hyperparameters}
\todonote{Auto-generated: search spaces from \texttt{models/classical.py} (and
quantum model files when they land); locked winning configurations from
\texttt{experiments/configs/e1\_locked/} (and successors); the complete
tuning-trial log lives in \texttt{results/e1\_tuning\_trials.csv} and is
summarized here.}

\chapter{Extended Result Tables}
\label{app:results}
\todonote{Per-metric master tables (accuracy, balanced accuracy, F1, ROC-AUC;
tuned and default variants) from analysis/output/, plus per-(seed, fold) raw
logs pointer.}

\chapter{Guide to the Code Repository}
\label{app:repo}
\todonote{Layout walkthrough (data/, models/, experiments/, analysis/,
results/, notebooks/); the reproduce-everything recipe (environment.yml,
\texttt{make\_splits --verify}, run.py configs, analysis scripts); the append-only
results contract and config-hash conventions; how the stored splits guarantee
comparability across the entire thesis.}


\bibliographystyle{abbrvnat}
\bibliography{refs}

\end{document}
